\section{Best Time to Buy and Sell Stock III}
\label{sec:dp-stock-3}

\problemstatement{Given $\textit{price}$, find the maximum profit from
\textbf{at most 2} transactions.}

\begin{patternbox}
\emph{``At most $2$ transactions''} requires a \textbf{new dimension}
added to Section~\ref{sec:dp-stocks-intro}'s template: a single
\textit{holding} bit can no longer distinguish ``this is my first
buy/sell cycle'' from ``this is my second,'' so we add a
\textit{transactionState} counter tracking exactly which of the four
actions (buy$1$, sell$1$, buy$2$, sell$2$) comes next.
\end{patternbox}

\subsection*{Deriving the recurrence}

Encode $\textit{cap} \in \{1,2,3,4\}$: $\textit{cap}=1$ means ``next
action is buy $1$,'' $2$ means ``next is sell $1$,'' $3$ means ``next is
buy $2$,'' $4$ means ``next is sell $2$.'' Odd \textit{cap} values are
buy-turns; even values are sell-turns:
\[
  f(i, \textit{cap}) =
  \begin{cases}
    \max\big(-\textit{price}[i] + f(i{+}1, \textit{cap}{+}1),\
    f(i{+}1, \textit{cap})\big) & \textit{cap} \text{ odd (buy turn)}, \\
    \max\big(\textit{price}[i] + f(i{+}1, \textit{cap}{+}1),\
    f(i{+}1, \textit{cap})\big) & \textit{cap} \text{ even (sell turn)}.
  \end{cases}
\]
Base case: $f(n, \cdot) = 0$, and also $f(i, 5) = 0$ for any $i$ (both
transactions already used up; no further action is possible, regardless
of how many days remain).

\subsection*{Approach 1: Recursion}

\begin{lstlisting}
#include <bits/stdc++.h>
using namespace std;

int stock3Recursive(int i, int cap, vector<int>& price) {
    int n = price.size();
    if (cap == 5 || i == n) return 0;

    if (cap % 2 == 1) { // buy turn
        return max(-price[i] + stock3Recursive(i + 1, cap + 1, price),
                   stock3Recursive(i + 1, cap, price));
    }
    return max(price[i] + stock3Recursive(i + 1, cap + 1, price),
               stock3Recursive(i + 1, cap, price));
}

int main() {
    vector<int> price = {3, 3, 5, 0, 0, 3, 1, 4};
    cout << stock3Recursive(0, 1, price) << endl; // 6
    return 0;
}
\end{lstlisting}

\begin{complexitybox}[Approach 1 Complexity]
\textbf{Time.} $O(2^n)$.
\textbf{Space.} $O(n)$ recursion stack.
\end{complexitybox}

\subsection*{Approach 2: Memoization}

\begin{lstlisting}
#include <bits/stdc++.h>
using namespace std;

int stock3Memo(int i, int cap, vector<int>& price, vector<vector<int>>& memo) {
    int n = price.size();
    if (cap == 5 || i == n) return 0;
    if (memo[i][cap] != -1) return memo[i][cap];

    int result;
    if (cap % 2 == 1) {
        result = max(-price[i] + stock3Memo(i + 1, cap + 1, price, memo),
                      stock3Memo(i + 1, cap, price, memo));
    } else {
        result = max(price[i] + stock3Memo(i + 1, cap + 1, price, memo),
                      stock3Memo(i + 1, cap, price, memo));
    }
    return memo[i][cap] = result;
}

int main() {
    vector<int> price = {3, 3, 5, 0, 0, 3, 1, 4};
    int n = price.size();
    vector<vector<int>> memo(n, vector<int>(6, -1));
    cout << stock3Memo(0, 1, price, memo) << endl; // 6
    return 0;
}
\end{lstlisting}

\begin{complexitybox}[Approach 2 Complexity]
\textbf{Time.} $O(n \cdot 4) = O(n)$.
\textbf{Space.} $O(n)$ memo $+$ $O(n)$ recursion stack.
\end{complexitybox}

\subsection*{Approach 3: Tabulation}

\begin{lstlisting}
#include <bits/stdc++.h>
using namespace std;

int stock3Tabulation(vector<int>& price) {
    int n = price.size();
    vector<vector<int>> dp(n + 1, vector<int>(6, 0));

    for (int i = n - 1; i >= 0; i--) {
        for (int cap = 4; cap >= 1; cap--) {
            if (cap % 2 == 1) {
                dp[i][cap] = max(-price[i] + dp[i + 1][cap + 1], dp[i + 1][cap]);
            } else {
                dp[i][cap] = max(price[i] + dp[i + 1][cap + 1], dp[i + 1][cap]);
            }
        }
    }
    return dp[0][1];
}

int main() {
    vector<int> price = {3, 3, 5, 0, 0, 3, 1, 4};
    cout << stock3Tabulation(price) << endl; // 6
    return 0;
}
\end{lstlisting}

\subsection*{Dry run}

\dryrun{Take $\textit{price} = [3,3,5,0,0,3,1,4]$.}

The optimal two transactions: buy at price $0$ (day $3$ or $4$), sell
at price $3$ (day $5$, profit $3$), buy at price $1$ (day $6$), sell at
price $4$ (day $7$, profit $3$): total $3+3=6$, matching
$\textit{dp}[0][1]=6$.

\begin{complexitybox}[Approach 3 Complexity]
\textbf{Time.} $O(n)$ ($4$ cap values per day).
\textbf{Space.} $O(n)$ table, no recursion stack.
\end{complexitybox}

\subsection*{Approach 4: Space Optimization}

\begin{lstlisting}
#include <bits/stdc++.h>
using namespace std;

int stock3SpaceOptimized(vector<int>& price) {
    int n = price.size();
    vector<int> next(6, 0), current(6, 0);

    for (int i = n - 1; i >= 0; i--) {
        for (int cap = 4; cap >= 1; cap--) {
            if (cap % 2 == 1) {
                current[cap] = max(-price[i] + next[cap + 1], next[cap]);
            } else {
                current[cap] = max(price[i] + next[cap + 1], next[cap]);
            }
        }
        next = current;
    }
    return next[1];
}

int main() {
    vector<int> price = {3, 3, 5, 0, 0, 3, 1, 4};
    cout << stock3SpaceOptimized(price) << endl; // 6
    return 0;
}
\end{lstlisting}

\begin{complexitybox}[Approach 4 Complexity]
\textbf{Time.} $O(n)$.
\textbf{Space.} $O(1)$ -- a fixed-size array of $6$ entries per day.
\end{complexitybox}

\begin{takeawaybox}
Capping the number of transactions adds exactly one new state dimension
(here, $4$ values, encoding ``which of the $4$ actions comes next''),
with the buy/sell branching logic itself unchanged from
Section~\ref{sec:dp-stocks-intro}'s template -- only which array slot
each branch writes into and reads from has shifted. This same
``add a counter dimension to cap a repeated action'' idea generalizes
directly to Stock IV's arbitrary $k$.
\end{takeawaybox}
